\section{Discussion}
\label{sec:discussion}

\textbf{What survives scrutiny.} Four results are robust to the controls we pre-specified. (i) Claim type is recoverable from transcribed driver speech well above a context-only baseline, on held-out races, the held-out season and unseen drivers. (ii) A compact instruction-tuned LLM with category definitions is the top claim classifier, needing no training. (iii) Vocal arousal can be extracted from audio, and partial fine-tuning of WavLM significantly enhances it compared to fixed features. (iv) The lap-time residual model is precise enough to function as an outcome. The measurement of timestamp lag is a useful insight for anyone entering this dataset for telemetry.

\textbf{What fails.} Speaker attribution misses its threshold. The structural explanation is that per-channel voice clusters need to be \emph{identified}, and the identification signal (linguistic hints) is feeble: text alone attributes speakers with only 0.39 accuracy. Prompt biasing has minimal impact on WER in compressed radio. No claim type is supported by telemetry beyond its base rate, and if any signal is present (\textsc{incident}), the reverse-time placebo indicates the event occurs before the message: it reflects the broadcaster's choice of significant moments, rather than the driver's account. This is precisely the confusion the placebo aimed to reveal, and we present it accordingly.

\textbf{Applying transfer learning with this data volume.} Three transfer strategies were compared on the same claim task. Full fine-tuning of a 184M-parameter encoder on $\sim$250 examples did not beat a linear model on lexical features, while an instruction-tuned LLM used zero-shot beat both by a wide margin. The difference is where the task knowledge comes from: the fine-tuned model must infer each category's boundary from few, imbalanced examples, whereas the LLM is \emph{told} the boundary in the same definitions the annotator used. For a taxonomy defined by an external criterion (what a sensor can check), stating that criterion beats learning it from a few hundred labels. Arousal behaves differently: it is expressed diffusely (intensifiers, profanity, prosody), and there fine-tuning paid off, both for text (DeBERTa was the best text arousal model) and especially for audio, where unfreezing WavLM's upper layers gave $+0.11$ AUC-PR. The learning curve supports this reading: DeBERTa's claim performance rises steeply with data but flattens below the LLM.

\textbf{Limitations.} (1) One annotator: no inter-annotator $\kappa$; label noise bounds every model. (2) Small samples: 317 driver utterances, five claim classes after dropping those under 30 (the brief required 100), and only 37 test-season items, so the chronological test is reported but not relied upon. (3) Attribution errors were removed by hand for labelled data only; unlabelled utterances used in T4 still contain $\sim$30\% non-driver speech, which attenuates any arousal effect. (4) The arousal annotator could read the transcript, which may inflate text-based arousal scores. (5) Telemetry lacks brake pressure, temperatures and steering, so several claims are uncheckable by design (\textsc{no\_sensor}). (6) Broadcast selection means radio is not a random sample of what drivers say.

\textbf{Reproducibility and ethics.} Models were implemented with Hugging Face Transformers~\cite{wolf2020transformers}, scikit-learn~\cite{pedregosa2011sklearn} and statsmodels~\cite{seabold2010statsmodels}. All numbers are produced by scripts in the project repository and written to \texttt{results/} as JSON/CSV; analysis choices were recorded before the labels they concern existed. Race audio is broadcast material: despite the dataset's CC~BY~4.0 tag, it is used here for coursework only and is not redistributed. The T2 target is \emph{telemetry-consistency}: a claim the telemetry does not corroborate is not thereby false.

\section{Conclusion and Future Work}
A carefully controlled pipeline on public F1 radio yields clear positives for the language task (claim types are identifiable from speech, best of all by a zero-shot LLM told the category definitions) and for paralinguistics (arousal is audible, and partial fine-tuning of a speech encoder helps substantially), alongside honest negatives for speaker attribution, ASR adaptation, telemetry corroboration and arousal's effect on pace. The negative controls are what make the positives credible. The most valuable next steps are a second annotator, more labels for rare claim types, a supervised attribution model trained on the hand labels, and LLM-assisted pre-labelling checked by humans.
